%==========================================================================================
\section{Follow-up (21 October, evening): why the original wins DA, and a multiscale decoder}\label{sec:round3}

Three questions were left open by \S\ref{sec:fairh8}. Every run in this section used the same DA
code, problems and tuning procedure as before. Results are in
\texttt{results/h8\_summary.json}.

\subsection{Was the curriculum model under-tuned?}

\textbf{No.} On the widened learning-rate grid its optimum stays at 1.0, inside the grid
(tuning error 0.022 / 0.019 / 0.021 / 0.031 at 0.3 / 1 / 3 / 10). Its DA numbers reproduced
exactly, so the values in \S\ref{sec:da} stand.

\subsection{Does checkpoint selection explain the original's DA advantage?}\label{sec:selection}

\textbf{Background.} Every DA result for your original KAE so far used
\texttt{best\_model.pth}. That checkpoint comes from epoch 40 and was picked by forecast loss
on \texttt{gt\_interp}, the same trajectories DA is tested on. The new models never see the
test trajectories before testing:
\begin{itemize}
\item the fair-horizon model uses its final checkpoint;
\item the curriculum model's checkpoint is picked on held-out Mach numbers.
\end{itemize}
I therefore re-ran the full DA suite on the original's \emph{final} checkpoint (epoch 199).
This is also the checkpoint behind its reproducible forecast numbers (11.43 / 2.12 / 18.62).
Its tuned learning rate is unchanged at 0.3, inside the grid.

\begin{center}\footnotesize
\begin{tabular}{lccccc}
\toprule
Sweep & \multicolumn{2}{c}{Original KAE} & Curriculum & H8, 256 & H8, 256\\
 & selected (ep.\ 40) & final (ep.\ 199) & & 3 seeds & multiscale\\
\midrule
\multicolumn{6}{l}{\emph{df = 9}}\\
window length $\delta_l$ & 0.085 & 0.099 & \textbf{0.094} & 0.096 & 0.109\\
noise                    & 0.078 & 0.101 & \textbf{0.094} & \textbf{0.094} & 0.105\\
sparsity                 & 0.086 & 0.104 & 0.100 & \textbf{0.098} & 0.112\\
misspecified noise       & 0.078 & 0.102 & 0.099 & \textbf{0.093} & 0.106\\
number of observations   & 0.079 & 0.106 & \textbf{0.095} & \textbf{0.095} & 0.105\\
joint sparsity           & 0.078$^\ast$ & 0.106 & 0.100 & \textbf{0.097} & 0.107\\
\midrule
\multicolumn{6}{l}{\emph{df = 1 (H8, 256: seed 0)}}\\
first-observation lead $\delta_f$ & 0.070 & 0.083 & \textbf{0.074} & 0.076 & 0.091\\
single observation       & 0.089 & 0.103 & 0.091 & \textbf{0.080} & 0.097\\
window length $\delta_l$ & 0.027 & 0.028 & \textbf{0.027} & 0.040 & 0.051\\
noise                    & 0.028 & \textbf{0.028} & \textbf{0.028} & 0.041 & 0.053\\
sparsity                 & 0.031 & \textbf{0.032} & 0.037 & 0.047 & 0.062\\
misspecified noise       & 0.028 & \textbf{0.029} & \textbf{0.029} & 0.041 & 0.053\\
number of observations   & 0.026 & 0.028 & \textbf{0.027} & 0.042 & 0.054\\
joint sparsity           & 0.027 & 0.029 & \textbf{0.028} & 0.043 & 0.057\\
\midrule
tuned learning rate      & 0.3 & 0.3 & 1.0 & 3.0 & 3.0\\
\bottomrule
\end{tabular}
\end{center}
{\footnotesize Bold: best among the models whose checkpoint never saw the test trajectories (all
columns except the first). $^\ast$ 3 of 5 points.}

\textbf{Result: yes, at df = 9 the advantage comes from checkpoint selection.}
\begin{itemize}
\item \textbf{df = 9:} the original's final checkpoint scores 0.099--0.106, \emph{worse} than
  both new models. The fair-horizon model (3 seeds) is 3--10\% better than it on every sweep.
\item \textbf{df = 1:} selection barely matters on the window sweeps (0.026--0.031 vs
  0.028--0.032). It does matter on the $\delta_f$ and single-observation sweeps (0.070 / 0.089
  vs 0.083 / 0.103). On the window sweeps the original and the curriculum model are level and
  best. The fair-horizon model is about
  1.5$\times$ worse on the window sweeps, but better on the single-observation sweep.
\item \textbf{The fair comparison is therefore:} the new model is better at df = 9 and
  level or worse at df = 1, depending on the variant.
\end{itemize}

\subsection{Why: how well do the observations pin down the initial state?}\label{sec:conditioning}

\textbf{Method} (\texttt{da\_conditioning.py}). On the df = 9 default problem, each model is
linearised around its best representation of the true state. Two Jacobians are compared:
\begin{itemize}
\item the map from latent to observations (what 4D-Var fits);
\item the map from latent to the initial state (what the analysis is scored on).
\end{itemize}
Their generalised singular values (the ``gains'') measure, in each direction, how much the
observations change per unit change of the initial state. The measure does not depend on how
the latent is parametrised, so models with different latents are comparable. A small gain
marks a part of the initial state that the observations barely see. The Jacobians use finite
differences and agree with autograd to $<10^{-3}$. Averages are over 8 problems.

\begin{center}\footnotesize
\begin{tabular}{lcccccc}
\toprule
 & \multicolumn{2}{c}{Original KAE} & Curriculum & H8, 256 & H8, 128 & H8, 256\\
 & ep.\ 40 & ep.\ 199 & & & & multiscale\\
\midrule
analysis error (this problem)              & 0.079 & 0.107 & 0.102 & 0.093 & 0.111 & 0.112\\
weakest gain                               & 0.10 & \textbf{0.02} & 0.13 & 0.11 & 0.15 & 0.11\\
10th-percentile gain                       & 0.44 & \textbf{0.14} & 0.45 & 0.43 & 0.42 & 0.45\\
directions with gain $<$ 0.1 $\times$ median & 7.5\% & \textbf{13.4\%} & 6.8\% & 6.6\% & 4.6\% & 6.4\%\\
modes decaying before the first observation & 0.8\% & \textbf{15.6\%} & 0 & 0 & 0 & 0\\
condition number in latent coordinates     & $3\times10^5$ & $4\times10^7$ & $2\times10^3$ & $2\times10^4$ & $8\times10^3$ & $1\times10^4$\\
misfit at analysis / at encoded truth      & 0.20 & 0.24 & 0.55 & 0.83 & 0.84 & 0.88\\
\bottomrule
\end{tabular}
\end{center}

\textbf{Result.}
\begin{itemize}
\item \textbf{It explains the two original checkpoints.} In the final checkpoint the weakly
  observed directions are three times weaker (10th-percentile gain 0.14 vs 0.44), and there
  are nearly twice as many of them. 16\% of its modes decay before the first observation
  (frame 9). That information is lost before 4D-Var can see it, which is why it
  is worse at df = 9 and not at df = 1. The selected epoch-40 checkpoint happens to lack
  these modes.
\item \textbf{It does not separate the selected original from the new models.} Their gain
  statistics are similar, and every new model has no fast-decaying modes and a far better
  latent condition number. What remains between them at df = 9 is not visible in these
  linear measures.
\item \textbf{Every model fits the observations better than its own encoded truth does} (last
  row $<1$). So the analysis error is limited by model error, not by optimisation, and more
  iterations would not help.
\end{itemize}

\subsection{Multiscale decoder}\label{sec:multiscale}

\textbf{Design} (\texttt{z\_inject} in \texttt{models/ct\_kae.py};
\texttt{configs/ctkae/fair\_h8\_ms.yaml}).
\begin{itemize}
\item The latent $z$ also writes 16 / 4 / 2 channels directly onto the decoder's 16$\times$32 /
  32$\times$64 / 64$\times$128 grids. This is a linear, POD-like path to the fine scales,
  motivated by POD's lower reconstruction error.
\item The new path enters through zero-initialised layers, so training starts from the plain
  decoder.
\item It is still a function of $(z, Ma)$ only, so the dynamics stay exactly linear.
\item It is trained end-to-end in the same 2 $\to$ 8 setting.
\item Cost: 33.1\,M parameters instead of 24.7\,M, and 3\% more time per step.
\end{itemize}
Rule fixed before the run: seeds 1--2 would follow only if seed 0's first-step error fell by at
least 10\% (below 0.235).

\begin{center}\small
\begin{tabular}{lccccccc}
\toprule
Seed 0, test & interp & extrap & long & lead 1 & lead 10 & energy at 58 & DA df = 9\\
\midrule
H8, 256              & \textbf{3.43} & \textbf{1.67} & 16.80 & 0.261 & 0.466 & 0.94 & \textbf{0.091--0.095}\\
H8, 256, multiscale  & 3.67 & 1.97 & \textbf{16.26} & 0.260 & \textbf{0.427} & 0.94 & 0.105--0.112\\
\bottomrule
\end{tabular}
\end{center}

\textbf{Result: no gain.}
\begin{itemize}
\item The first-step error is unchanged (0.260 vs 0.261), so the rule stopped the extra
  seeds.
\item Interpolation and extrapolation are at the top of the plain model's seed spread
  (3.38--3.66 and 1.67--2.29).
\item DA is worse.
\item Zero-shot time stepping still holds (3.683 at half steps vs 3.672 at whole steps).
\item \textbf{What this rules out.} Together with the latent-size test (\S\ref{sec:repr}),
  this rules out decoder capacity and latent size as the cause of the first-step error. The
  same decoder reaches 0.08 under one-step training (\S\ref{sec:acdmsetting}). The floor is
  therefore set by the trade-off in the training objective between reconstruction and
  multi-step prediction.
\end{itemize}
